% ============================================
% Johns Hopkins Letter Template
% Assistant Professor Cover Letter – University of California, Merced
% ============================================

\documentclass[11pt,letterpaper]{letter}

\usepackage{graphicx}
\usepackage{eso-pic}
\usepackage{geometry}
\usepackage{microtype}
\usepackage[hidelinks]{hyperref}

% ----- Page Setup -----
\geometry{left=25mm,right=25mm,top=25mm,bottom=25mm}

% ----- Logo Placement (foreground, first page only) -----
\newcommand{\PutJHULogoFG}[1]{%
  \AddToShipoutPictureFG*{%
    \AtPageUpperLeft{%
      \hspace{10mm}%
      \raisebox{-38mm}{%
        \includegraphics[width=6cm]{#1}%
      }%
    }%
  }%
}

% ----- Sender Information -----
\signature{Decory Edwards}
\address{Department of Economics \\ Johns Hopkins University \\ Baltimore, MD 21218 \\ \texttt{dedwar65@jh.edu}}

\begin{document}
\PutJHULogoFG{Images/jhulogo.png}

\begin{letter}{Search Committee \\
Department of Economics \& Business Management \\
University of California, Merced \\
5200 North Lake Road \\
Merced, CA 95343 \\ USA}

\opening{Dear Members of the Search Committee,}

I am writing to express my interest in the tenure‐track Assistant Professor position in Economics at the University of California, Merced, beginning July 1, 2026. I am a Ph.D. candidate in Economics at Johns Hopkins University, advised by Professors Christopher D. Carroll, Nicholas W. Papageorge, and Stelios Fourakis, and I expect to receive my degree in May 2026. My research fields are macroeconomics, wealth and income dynamics, and computational economics.

My research agenda is unified by a focus on understanding the determinants of wealth inequality and the mechanisms through which heterogeneity in returns to wealth shapes aggregate outcomes. In my job‐market paper, I estimate the degree of heterogeneity in individual portfolio returns using micro‐data from the Survey of Consumer Finances and simulate a structural model in which households face idiosyncratic return risk. I show that allowing for persistent heterogeneity in returns substantially improves the model’s ability to match the empirical distribution of wealth, offering a tractable way to connect micro‐level return dispersion to macroeconomic inequality.

In a second project, I use the Health and Retirement Study to examine the relationship between interpersonal trust and portfolio performance. Building on the literature that finds a hump‐shaped relationship between trust and income, I show that a similar relationship holds for individual returns to wealth measured in the HRS data, highlighting a behavioral channel through which social attitudes shape saving and investment outcomes. This work also provides new estimates of the persistent component of returns to net worth, which motivate the calibration of heterogeneity in my structural model.

The University of California, Merced’s Department of Economics \& Business Management is forward‐looking, engaged with both rigorous research and inclusive teaching, and situated in a region that values community impact and interdisciplinary collaboration. I am attracted to the department’s commitment to combining academic excellence with access and innovation, and I believe my work on wealth dynamics and computational macroeconomics complements the department’s mission and strengths.

\textbf{Teaching.} My teaching experience centers on undergraduate macroeconomics and an upper‐level macro‐policy seminar, where I have led weekly discussion sections and maintained regular office hours for classes ranging from 16 to 40 students. My approach is student‐centered: I aim to help students “convince themselves” that they have moved from confusion to understanding by providing structured office‐hour coaching that builds trust and confidence. At UC Merced, I am prepared to teach \textit{Principles of Macroeconomics}, \textit{Intermediate Macroeconomic Theory}, \textit{Money and Banking}, and \textit{Quantitative Methods for Economics}. I would also welcome developing electives such as \textit{Economics of Inequality and Tax Policy} or \textit{Computational Macroeconomics and Policy}, emphasizing reproducible workflows and evidence‐based decision‐making.

I believe I would be a strong fit because my computational and empirical toolkit allows me to (i) produce robust, data‐backed estimates of returns and distributional outcomes, (ii) build and validate models that speak to macroeconomic and policy‐relevant questions, and (iii) collaborate with colleagues and engage undergraduate and graduate students in research and teaching. I am enthusiastic about the opportunity to join UC Merced’s community and contribute to the Department’s teaching, research, and service missions.

I would be honored to join the faculty at the University of California, Merced and to mentor students pursuing quantitative careers in economics, business, and policy. Thank you for your consideration; I welcome the opportunity to discuss my fit with your department at your convenience.

\closing{Sincerely,}

\end{letter}
\end{document}
